% 同一可行多边形上两个目标方向：(a) max 3x1+2x2 只在顶点(4/3,4/3)达到20/3；
% (b) max 2x1+x2 的等值线与边 2x1+x2=4 平行，整条边[(2,0),(4/3,4/3)]都最优，最优值4。
\begin{tikzpicture}[x=10mm,y=10mm,
  note/.style={font=\FigureNoteFont\fontsize{7.5}{9}\selectfont,text=FigureMuted},
  ttl/.style={font=\fontsize{8}{10}\selectfont,anchor=north}]
  \begin{scope}
    \fill[FigureBlueFill] (0,0)--(2,0)--(4/3,4/3)--(0,2)--cycle;
    \draw[FigureBlue,line width=.9pt] (0,0)--(2,0)--(4/3,4/3)--(0,2)--cycle;
    \begin{scope}\clip (0,0) rectangle (3.2,3.6);
      \foreach \k in {2,4}{\draw[FigureGreen,line width=.6pt,dashed] (0,\k/2)--(\k/3,0);}
      \draw[FigureGreen,line width=1.1pt] (0,10/3)--(20/9,0);
    \end{scope}
    \draw[->,FigureStroke,line width=.6pt] (-0.15,0)--(3.2,0) node[right]{$x_1$};
    \draw[->,FigureStroke,line width=.6pt] (0,-0.15)--(0,3.6) node[above]{$x_2$};
    \draw[->,FigureRed,line width=1pt] (4/3,4/3)--++(0.67,0.45) node[right]{$(3,2)$};
    \fill[FigureRed] (4/3,4/3) circle (1.7pt);
    \node[ttl] at (1.5,-0.35) {(a) 唯一最优顶点};
  \end{scope}
  \begin{scope}[xshift=50mm]
    \fill[FigureBlueFill] (0,0)--(2,0)--(4/3,4/3)--(0,2)--cycle;
    \draw[FigureBlue,line width=.9pt] (0,0)--(2,0)--(4/3,4/3)--(0,2)--cycle;
    \begin{scope}\clip (0,0) rectangle (3.2,3.6);
      \foreach \k in {1,2,3}{\draw[FigureGreen,line width=.6pt,dashed] (0,\k)--(\k/2,0);}
      \draw[FigureGreen,line width=.6pt] (0,4)--(2,0);
    \end{scope}
    \draw[FigureRed,line width=2pt] (2,0)--(4/3,4/3);
    \draw[->,FigureStroke,line width=.6pt] (-0.15,0)--(3.2,0) node[right]{$x_1$};
    \draw[->,FigureStroke,line width=.6pt] (0,-0.15)--(0,3.6) node[above]{$x_2$};
    \draw[->,FigureRed,line width=1pt] (5/3,2/3)--++(0.72,0.36) node[right]{$(2,1)$};
    \fill[FigureRed] (2,0) circle (1.7pt);
    \fill[FigureRed] (4/3,4/3) circle (1.7pt);
    \draw[FigureRed,fill=white,line width=.8pt] (5/3,2/3) circle (1.5pt);
    \node[note,anchor=west] at (1.85,1.55) {最优面：整条边};
    \node[ttl] at (1.5,-0.35) {(b) 最优面是一条边};
  \end{scope}
\end{tikzpicture}
